\documentclass[a4paper,12pt]{article}
\usepackage[utf8]{inputenc}
\usepackage{graphicx}
\usepackage{hyperref}
\usepackage{listings}
\usepackage{xcolor}
\usepackage{geometry}

\geometry{
 a4paper,
 total={170mm,257mm},
 left=20mm,
 top=20mm,
}

\title{\textbf{NeuroGraph: A Self-Improving Agentic Architecture for Graph Retrieval-Augmented Generation}}
\author{Avirup \\ \textit{Department of Computer Science}}
\date{\today}

\begin{document}

\maketitle

\begin{abstract}
Retrieval-Augmented Generation (RAG) has emerged as a powerful technique to ground Large Language Models (LLMs) in external data. However, traditional vector-based RAG often lacks structured reasoning and a mechanism for long-term self-correction. In this paper, we introduce \textbf{NeuroGraph}, an experimental agentic architecture that combines GraphRAG with an automated Reinforcement Learning from Human Feedback (RLHF) loop. The system employs a tri-agent approach---The Builder, The Actor, and The Critic---to dynamically construct a Knowledge Graph from user interactions while simultaneously creating a dataset of failures for future fine-tuning. This approach allows the system to not only retrieve context but to actively learn from user dissatisfaction in real-time.
\end{abstract}

\section{Introduction}
Large Language Models (LLMs) often suffer from "hallucinations" and a lack of long-term consistency. While RAG systems address the knowledge gap by retrieving documents, they are typically static; they do not learn or adapt based on their mistakes during interaction. Furthermore, vector databases capture semantic similarity but struggle with complex structural relationships between entities.

To address these limitations, NeuroGraph implements two key innovations:
\begin{enumerate}
    \item \textbf{Graph-Based Memory:} Using a Graph Database (Neo4j) to map entities and relationships, providing structured context.
    \item \textbf{The "Critic" Loop:} A background agent that monitors user sentiment to identify and log failures, enabling an automated data fly-wheel for Direct Preference Optimization (DPO).
\end{enumerate}

\section{System Architecture}
The NeuroGraph architecture consists of three specialized agents working in tandem.

\subsection{1. The Builder (Graph Manager)}
The Builder is responsible for long-term memory acquisition. It interfaces with the Neo4j database to store information as triplets: $(Subject) \xrightarrow{RELATION} (Object)$.
\begin{itemize}
    \item \textbf{Extraction:} The agent processes user input (e.g., "I am building a React app") and utilizes an LLM to extract structured data: \texttt{\{source: 'User', relation: 'BUILDING', target: 'React App'\}}.
    \item \textbf{Merging:} These triplets are merged into the graph, ensuring no duplication while building a dense web of knowledge over time.
\end{itemize}

\subsection{2. The Actor (RAG Engine)}
The Actor serves as the primary interface for user interaction.
\begin{itemize}
    \item \textbf{Context Retrieval:} Unlike simple vector search, the Actor queries the graph for direct neighbors of the entities mentioned in the user's prompt. This retrieves explicit relationships (e.g., \textit{User LIKES Python}) to ground the generation.
    \item \textbf{Generation:} The retrieved graph context is injected into the system prompt of a local LLM (Llama 3.1) to generate a factually grounded response.
\end{itemize}

\subsection{3. The Critic (Dissatisfaction Monitor)}
The Critic addresses the challenge of self-improvement. It runs asynchronously to analyze the quality of the interaction.
\begin{itemize}
    \item \textbf{Sentiment Analysis:} The Critic evaluates the user's latest response to the previous bot reply. It looks for signals of frustration (e.g., "No, that's wrong", "Stop").
    \item \textbf{Data Collection:} Upon detecting dissatisfaction, the Critic logs the interaction pair (User Input, Bot Reply, Reason) into a \texttt{dissatisfaction\_dataset.jsonl}. This synthetic dataset serves as the ground truth for future alignment training (DPO or Fine-Tuning).
\end{itemize}

\section{Implementation Details}
The prototype is implemented in Python, leveraging the following technologies:
\begin{itemize}
    \item \textbf{Database:} Neo4j Community Edition for graph storage.
    \item \textbf{LLM Backend:} Ollama running Llama 3.1 (8B Parameters).
    \item \textbf{Orchestration:} LangChain for prompt management and chain execution.
\end{itemize}

\section{Conclusion and Future Work}
NeuroGraph demonstrates a viable path toward self-correcting AI systems. By closing the loop between retrieval and user feedback, we move from static information retrieval to dynamic, learning agents. Future work involves automating the fine-tuning pipeline, where the logged dissatisfaction dataset automatically triggers a DPO training run to update the model weights.

\end{document}
